\SourceLeaf{178}seems to us more in keeping with the dignity of science than the vain
endeavor to conceal the gap that cannot be filled under an untenable web
of sham proofs.

The author of the first work had already made a similar attempt 15 years earlier in a short
treatise, »Tentamen novae parallelarum theoriae notione situs fundatae«.
He sought to base everything on the concept of identity
of position. He defines parallel lines as those
straight lines that have the same position and concludes from this that such lines
must be intersected by any third straight line at equal
angles, since these angles are nothing other
than the measure of the difference between the position of this third line and the
positions of the two parallel lines. This form of proof is repeated in the present new
work, although we cannot say that it has gained in strength through the
philosophical considerations woven into it. Every geometer must
dispute the assertion on p. 24, »Notionem situs e geometria adeo non excludi posse, ut potius
notionibus ejus fundamentalibus annumeranda sit, \textls{dudum omnes agnovere
geometrae}«, in the sense in which the author uses the concept of position in his
proof. If we begin with the
author's definition, »Situs est modus, quo plura coëxistunt vel juxta se
existunt in spatio«, then position is a purely relational concept, and we
may say that two straight lines $A,B$ have a certain position relative to one another,
which is the same as the mutual position of two other lines $C,D$. But the
author uses the word position in his proof as an absolute concept, speaking
of identity of position between two noncoincident straight lines. This
meaning plainly remains empty and without foundation until we know what
we are to understand by such an identity and how we are to recognize
it. If it is to be recognized by the equality of angles with \textls{one} third
straight line, then without a prior proof we do
not yet know whether the very same equality also holds for the angles with a fourth straight
line. If equality of the angles with \textls{every} other straight line is to be the criterion, we again do not know whether
the same position is possible without coincidence. We accordingly stand
\textls{after} the author's proof at exactly the same point at which we stood
\textls{before} it.

